\section{Related Work}\label{sec:related}

\paragraph{Bugs in ZK circuits and proof systems}
Chaliasos et al.\ systematize vulnerabilities in SNARK systems and
applications~\cite{sok-snark-bugs}. Their taxonomy is organized by the layer
where a bug lives (circuit, front end, back end, integration), and
under-constrained circuits dominate it. Tools follow this focus.
circomspect~\cite{circomspect} lints Circom for dangerous constructs,
Picus~\cite{picus} uses SMT solving to decide whether a circuit uniquely
determines its outputs, ZKAP~\cite{zkap} detects circuit anti-patterns
statically, and zkFuzz~\cite{zkfuzz} fuzzes circuits for under- and
over-constraints. All of these analyze the relation $R$. A binding flaw can
exist with a perfectly constrained circuit, because it concerns which consumer
values the relation is \emph{meant} to fix. We complement this line of work
by characterizing and detecting the integration-level class that circuit
tools cannot see.
\TODO{W4}{literature search: ZK integration/verifier-misuse papers 2024--2026, zkVM integration bugs, recent circuit fuzzers; verify every citation}

\paragraph{Smart-contract analysis}
Static analyzers such as Slither~\cite{slither}, symbolic executors such as
Mythril~\cite{mythril}, and fuzzers such as Echidna~\cite{echidna} and Foundry
find reentrancy, access-control and arithmetic flaws in contracts. They can
follow taint from calldata to a transfer, but they have no notion of a proved
statement, so an unbound recipient looks like an ordinary parameter. Our
evaluation confirms this: with generic properties they find none of the
binding flaws, and a fuzzer becomes effective only once it is given the
binding invariant.

\paragraph{LLM-based vulnerability detection}
LLMs have been combined with static analysis to detect logic flaws in smart
contracts~\cite{gptscan} and used as agents that explore code with
tools~\cite{llm4vuln}. These systems rely on general vulnerability knowledge.
\tool{} differs in two ways. Its procedure is derived from a precise
definition and forces exhaustive, checkable ledgers instead of free-form
review, and it upgrades warnings to executable evidence.

\paragraph{Cross-chain verification}
Bridges and light clients also act on verified messages, and several large
bridge incidents involved a message being accepted without being bound to the
intended source, recipient or verifier~\cite{bridge-sok}. ZK bridges and zkVM
light clients inherit this risk at the proof-to-effect interface. Our P2 and
P4 patterns capture it, and several findings in our dataset come from such
systems.

\paragraph{Signature and cryptographic API misuse}
The patterns of Section~\ref{sec:definition} mirror long-studied misuse of
signatures and cryptographic APIs: unsigned fields, missing domain separation,
and nonce reuse~\cite{egele-crypto-misuse,cryptoguard}. Front-running, which
turns an unbound recipient into a theft, is well documented in
DeFi~\cite{flashboys}. To our knowledge, we are the first to define binding
between a ZK proof's statement and a contract's effect as a property, to
measure its prevalence, and to detect it automatically.
